\section{CPU、内存和设备通过互连交换信息}

CPU 需要地址来说明“访问哪里”，需要数据线传送内容，还需要读写和完成信号
说明事务类型与结果。内存和设备都可以接到互连上，但它们只响应分配给自己的
地址或端口。CPU 发出地址不等于数据已经到达，目标还可能要求等待或返回错误。

\begin{pdfdiagram}
  \node[os state] (cpu) {CPU\\地址、数据、读写};
  \node[os block,right=of cpu] (fabric) {平台互连\\选择目标};
  \node[os hardware,above right=of fabric] (rom) {ROM};
  \node[os hardware,right=of fabric] (ram) {RAM};
  \node[os hardware,below right=of fabric] (device) {设备寄存器};
  \draw[os bus] (cpu) -- (fabric);
  \draw[os arrow] (fabric) -- (rom);
  \draw[os arrow] (fabric) -- (ram);
  \draw[os arrow] (fabric) -- (device);
\end{pdfdiagram}
\diagramnote{概念图：平台按地址或端口选择目标；具体 x86 地址稍后给出。}

\section{RAM 保存运行中的可改数据}

RAM 可以读取和写入，断电后不保证保留内容。栈、页表、内核对象和用户程序
运行时的数据都会进入 RAM。CPU 刚复位时不能假设 RAM 已经包含可执行内核，
因为这些字节尚未被装入，而且部分 RAM 还可能需要平台初始化。

ROM 在这里表示“复位后已经存在、CPU 能读取的非易失字节”。项目把最早代码
做成 128 KiB 镜像，并让 QEMU 把它映射到 x86 复位地址附近。ROM 不是一种
汇编语言，也不是 CPU 自动产生的程序；它仍然由我们编译、链接和检查。

\section{取指从程序计数器给出的地址开始}

处理器把程序计数器中的地址送到取指接口，平台选择存储目标并返回若干字节。
译码器把字节解释成当前指令，执行单元改变寄存器或内存，程序计数器随后指向
下一条指令。跳转指令会改写这个顺序。

假设一台教学 CPU 的 PC 为 \code{0x1000}，该地址存放一条长度为 3 字节的
普通指令。提交后 PC 变为
\[
  0x1000+3=0x1003.
\]
如果这条指令是相对跳转，还要把有符号位移加到 \code{0x1003}。x86 复位向量
使用的正是这种计算。

\section{复位把第一条取指固定下来}

CPU 内部有许多寄存器。若上电后它们取值任意，平台就无法知道第一条访问会去
哪里。复位把架构状态带到处理器手册规定的值，其中包括取指位置。平台必须在
那个位置提供有效指令，否则处理器会读到空值、错误数据或直接停在异常路径。

项目使用的 x86-64 CPU 从物理地址 \code{0xFFFFFFF0} 附近开始取指。这个地址
位于 4 GiB 边界下方 16 字节。为什么 64 位 CPU 仍保留这一旧式入口，需要先
认识最小 x86 寄存器和跳转指令，再沿 8086 的兼容历史解释。

\section{第一次实验只追踪四个检查点}

启动的第一遍只回答四个问题：

\begin{enumerate}[label=\arabic*.]
  \item ROM 镜像在复位地址是否真的含有跳转操作码；
  \item CPU 是否从该地址取回了相同字节；
  \item 跳转目标是否仍位于 ROM 映射范围；
  \item 入口是否建立栈并写出第一条串口日志。
\end{enumerate}

此时不需要先掌握缓存一致性、DRAM 训练、PCB 阻抗或复杂流水线。它们都是真实
问题，但不应阻挡第一条指令。附录保留这些硬件深化内容，读者可以在需要解释
某个现象时再回来。
